\chapter{Introduction}\label{s:intro}

\section{Motivation}

Scientific computing increasingly depends on complex, multi-step computational
workflows: data must be fetched, cleaned, analysed, and interpreted, often
across heterogeneous compute infrastructure. Workflow Management Systems (WMSs)
such as Pegasus~\cite{deelman2015pegasus}, Nextflow~\cite{di2017nextflow}, and
Snakemake~\cite{molder2021snakemake} have become indispensable in fields like
genomics, epidemiology, and climate science precisely because they automate
resource allocation, fault recovery, and execution ordering. Yet, despite their
maturity, these systems impose a heavy cognitive cost on scientists. To express
even a straightforward intent---``compute GC content for sequence
ATGCGATCG''---a researcher must first understand the WMS's execution model,
write a domain-specific language (DSL) script that explicitly enumerates every
computational step, declare data dependencies as a directed acyclic graph (DAG),
and then debug execution failures at the infrastructure level rather than the
scientific level.

This mismatch between \emph{scientific intent} and \emph{workflow specification}
has been recognised as a fundamental barrier to reproducible science. Scientists
spend substantial effort debugging YAML configurations and DAG topologies instead
of advancing research questions. The cognitive overhead compounds further in
distributed settings, where data locality, containerisation, and security
policies add additional layers of complexity to what should be a simple
computational request.

Recent advances in Large Language Models (LLMs) open a new possibility: a
researcher could describe a computational task in plain language, and an
intelligent system could automatically translate that natural-language intent
into a runnable workflow. This vision---\emph{intent-based orchestration}---has
attracted growing interest in the systems and HPC communities. However, existing
approaches remain incomplete. Systems that accept natural language either
generate generic code that cannot be executed in a target WMS environment, or
they lack the formal provenance guarantees required for reproducible science.
Conversely, WMSs that provide strong reproducibility have no mechanism for
accepting high-level declarative intent.

The gap is thus precise and actionable: no system currently bridges declarative,
natural-language intent expression with the rigorous, provenance-tracked
execution guarantees of a production WMS.

\section{Problem Statement}

This thesis addresses the following core problem: given a scientific computing
infrastructure based on Brane---a distributed WMS designed for federated,
privacy-sensitive data analysis---how can a researcher specify computational
tasks in natural language and reliably obtain correct, reproducible workflow
executions?

The problem decomposes into several interacting challenges. First, natural
language is inherently ambiguous: the same scientific intent can be expressed
in many ways, and different phrasings should produce semantically equivalent
workflows. Second, Brane workflows are expressed in BraneScript, a typed DSL
with strict package invocation semantics; translation errors produce workflows
that either fail to compile or silently return incorrect results. Third,
scientific workflows require reproducibility: the same intent must produce the
same execution trace across reruns and across different sessions.

Addressing these challenges requires not only a translation mechanism from
natural language to BraneScript, but also a retrieval mechanism to ground
translation in available package capabilities, and a caching mechanism to
avoid redundant LLM calls for repeated or semantically equivalent intents.

\section{Research Questions}

This thesis investigates the following research questions:

\begin{enumerate}
    \item[\textbf{RQ1}] To what extent can LLMs translate natural-language
        scientific intents into syntactically and semantically correct
        BraneScript workflows?
    \item[\textbf{RQ2}] How does Retrieval-Augmented Generation (RAG) with
        package-aware context affect translation accuracy compared to a
        zero-shot baseline?
    \item[\textbf{RQ3}] How does supervised fine-tuning (SFT) on domain-specific
        BraneScript examples affect translation accuracy and generalisation?
    \item[\textbf{RQ4}] Can semantic caching effectively serve as a
        reproducibility layer for intent-based orchestration, and what are its
        fundamental limitations for scientific computing?
\end{enumerate}

\section{Contributions}

This thesis makes the following contributions:

\begin{enumerate}
    \item \textbf{NLP-to-BraneScript Translator.} A complete pipeline that
        accepts a natural-language scientific intent and produces a runnable
        BraneScript workflow via an LLM, with RAG-based package context retrieval
        using ChromaDB.

    \item \textbf{Curated BraneScript Dataset.} A dataset of 584 unique
        intent--BraneScript pairs spanning seven scientific domains (healthcare,
        genomics, epidemiology, statistics, text analysis, data masking, and
        datetime), constructed with explicit deduplication and diversity criteria.

    \item \textbf{Fine-Tuned Models.} Supervised fine-tuning experiments on
        Qwen3.5-9B and Qwen3.6-27B using LoRA, evaluated against base model
        baselines on BraneScript compilation and execution correctness.

    \item \textbf{Semantic Cache for Reproducible Orchestration.} A
        ChromaDB-backed semantic cache that maps intent embeddings to previously
        generated BraneScript workflows, enabling sub-millisecond response for
        repeated or paraphrased intents, with a benchmarked threshold sweep for
        precision--recall trade-off analysis.

    \item \textbf{Empirical Evaluation.} A systematic evaluation of model size,
        RAG depth, and caching threshold, yielding practical guidance on deploying
        LLM-based intent translators in production scientific computing
        environments.
\end{enumerate}

\section{Thesis Structure}

The remainder of this thesis is organised as follows. Chapter~\ref{s:background}
provides background on scientific workflow systems, provenance, intent-based
orchestration, and LLMs in scientific computing. Chapter~\ref{s:overview}
presents a high-level architecture of the proposed system. Chapter~\ref{s:design}
describes the design decisions. Chapter~\ref{s:implementation} details the
implementation. Chapter~\ref{s:evaluation} reports the empirical results.
Chapter~\ref{s:discussion} interprets the results in the context of the research
questions. Chapter~\ref{sec:threats} discusses threats to validity.
Chapter~\ref{s:related} surveys related work, and Chapter~\ref{s:conclusion}
concludes with future directions.




